\scanpage{38}
\mspage{4}
(15) étant définis donc directement en termes du 1-complexe \inserted{comb.} $K$ et de la donnée des trois ps. droites $D_i\subset K$ dans $K$. On peut réécrire (15) comme un syst. d'iso
\[
 \underbrace{\left(\bigwedge_{i\in I}\Omega_i\right)\wedge\operatorname{or}(I)}
 _{\substack{\text{or. de l'oct.}\\\text{défini par}\\(\Omega_i)_{i\in I}}}
 \simeq
 \underbrace{\left(\bigwedge_{i\in I}\Phi_i\right)\wedge\operatorname{or}(I)}
 _{\substack{\text{or. de l'oct.}\\\text{défini par}\\(\Phi_i)_{i\in I}}}
 \simeq\bigwedge_{i\in I}A_i\tag*{(15 bis)}
\]
Ceci posé, il y a, \deleted{une iso. can. entre} \inserted{du fait que} \uline{$K$ provient de la carte} $C$, \uline{une} \deleted{\unsure{iso}} \inserted{\uline{trivialisation}} \uline{can.} \uline{des} $\mathbb Z/2\mathbb Z$-\uline{torseurs} (15 bis). De façon précise, le foncteur
\diagram{(16)\quad
\begin{tikzcd}[ampersand replacement=\&,column sep=small]
\text{(oct. gauches)}\arrow[r]\&\substack{
\overbrace{\text{graphes octaédraux gauches}}^{\text{gr. oct. gauches}^{+}}\\
\left\{\begin{gathered}
\text{munis d'une trivialisation}\\
\bigwedge_{i\in I} A_i\simeq\mathbb{1}\quad\text{(où }I\text{ est l'ens}\\
\text{des \inserted{3} sommets, et }A_i\text{ est l'}\\
\text{ens des \inserted{2} arêtes qui ne passent}\\
\text{pas par le sommet }i
\end{gathered}\right.}
\end{tikzcd}}
\margin{on appellera cette structure suppl. une \uline{orientation} du gr. octaédral gauche}
est une \uline{équivalence} de catégories

On peut résumer les relations entre octaèdres et graphes octaédraux gauches ou ordinaires, orientés ou non, dans le diagramme de foncteurs suivant (entre groupoïdes connexes)
\diagram{(17)\quad
\begin{tikzcd}[ampersand replacement=\&,column sep=small,row sep=small]
\text{oct}^{+}\arrow[rr,"\approx"]\arrow[dr,"\approx"]\arrow[dd,hook] \& \& \text{oct. gauches}\arrow[dr,"\approx"]\arrow[dd,"\parallel" description] \& \\
\& \text{gr. oct.}^{+}\arrow[rr,"\approx",crossing over]\arrow[dd,hook] \& \& \text{gr oct. gauches}^{+}\arrow[dd,hook]\\
\text{oct.}\arrow[rr]\arrow[dr,"\approx"] \& \& \text{oct. gauche}\arrow[dr,""] \& \\
\& \text{gr. oct.}\arrow[rr,"\approx"] \& \& \text{gr. oct. gauche}
\end{tikzcd}}, 
% archive sheet number: 37
